\documentclass[xcolor=table, aspectratio=1610]{beamer}

\usetheme[]{Montpellier}

\setbeamertemplate{itemize items}[square]
\setbeamertemplate{section in toc}[sections numbered]

\setbeamertemplate{frametitle}{
    \begin{centering}
        \small\textsc{\insertframetitle}\par
    \end{centering}
}
\definecolor{mygreen}{RGB}{0, 128, 128}
\definecolor{mylightgreen}{RGB}{144,238,144}

\setbeamercolor*{palette primary}{bg=mygreen!70, fg=white}
\setbeamercolor*{palette secondary}{bg=mygreen!1, fg=black}
\setbeamercolor*{palette tertiary}{bg=mygreen!1, fg=black}
\setbeamercolor*{palette quaternary}{bg=mygreen!85, fg=white}

\setbeamercolor{block title}{bg=mygreen, fg=white}
\setbeamercolor{block body}{bg=mylightgreen, fg=black}
\setbeamertemplate{title page}[default][rounded=false]
\setbeamercolor{item}{fg=mygreen}
\setbeamercolor{frametitle right}{bg=mygreen!60}
\usepackage{amsfonts}
\usepackage{amsmath}
\usepackage{amsthm}
\usepackage{amssymb}
\usepackage{mathrsfs}
\usepackage{tikz}
\usepackage{graphicx}
\usepackage{subcaption}
\usepackage{booktabs}
\usepackage{makecell}
\usepackage{multirow}
\usepackage{multicol}
\usepackage{fouriernc}
\hypersetup{
    colorlinks=true,
    linkcolor=black,
    urlcolor=blue,
    citecolor=blue
}
\usetikzlibrary{arrows,calc}
\usepackage{relsize}

\makeatletter
\usepackage{xcolor}
\definecolor{customtitlebg}{RGB}{230, 230, 250}
\definecolor{customauthorbg}{RGB}{245, 245, 220}
\definecolor{customdatebg}{RGB}{240, 255, 240}

\setbeamercolor{title}{bg=mygreen, fg=white}

\makeatletter
\defbeamertemplate*{title page}{custom}{
  \vbox{}
  \vfill
  \begin{centering}
    \begin{beamercolorbox}[sep=14pt,rounded=false, shadow=true, wd=\textwidth, center]{title}
      \usebeamerfont{title}\inserttitle\par%
      \ifx\insertsubtitle\@empty%
      \else%
        \vskip0.25em%
        {\usebeamerfont{subtitle}\usebeamercolor[fg]{subtitle}\insertsubtitle}%
        \par%
      \fi%
    \end{beamercolorbox}%
    \vskip1em\par%
    \begin{beamercolorbox}[sep=8pt,center]{author}
      \usebeamerfont{author}\insertauthor\\
      {
\vspace{80pt}
\vspace{\fill}
\footnotesize{Last Updated: \today}
}
    \end{beamercolorbox}
    \vskip1em\par%
    \vfill
  \end{centering}
}
\makeatother
\useinnertheme{rounded}

\setbeamertemplate{title page}[custom]

\title[Ghana--Summary]{\textsf{Summary Data on Acid Cropland and Lime Requirements in Ghana}} 
\author[GAIA]{GAIA\\ \href{www.acidsoils.africa}{www.acidsoils.africa} } 
\institute[CIMMYT] 

\begin{document}

\begin{frame}[plain]
\textsc{\titlepage}
\end{frame}

\begin{frame}
\frametitle{Outlook}
\tableofcontents 
\end{frame}

\section{Introduction}
\begin{frame}{Introduction}
    \begin{itemize}
        \item The slide deck provides a \textbf{data-driven}, \textbf{ex-ante analysis} of soil acidity across \textbf{Ghana}, evaluating the scale of affected areas, potential crop yield impacts, and the expected benefits of lime application for proactive soil remediation.
        \bigskip
        \item \textbf{Highlights}:
        \begin{itemize}
            \item \textbf{Acidic Cropland}: Identifies regions with high acidity levels impacting crop yields.
            \item \textbf{Yield Loss}: Quantifies productivity losses in key crops due to acidic soils.
            \item \textbf{Lime Requirements}: Estimates the lime amounts needed to improve soil health in these areas.
            \item \textbf{Economic Potential}: Highlights estimated increases in production and production value from lime application.
        \end{itemize}
        \bigskip
        \item Provide policymakers, farmers, and stakeholders with actionable insights to prioritize areas with high potential for productivity gains and economic return from soil remediation.
    \end{itemize}
\end{frame}
\section{Acidic Cropland}
\subsection*{Top 10 Regions in Ghana with the highest acidic cropland}

\begin{frame}{Top 10 Regions in Ghana with the highest acidic cropland}
    \begin{figure}
    \includegraphics[width=0.95\textwidth]{D:/Dropbox/CIMMYT-WORK/Analysis/bisrat-analysis-workflows/gaia_exante_bisrat/gaia-ex-ante-wow-shared/data-output//figures_and_tables/GHA//GHA_plt_acidic_cropland.png}
    \end{figure}
\end{frame}

\subsection*{Top 10 Regions in Ghana with the highest acidic cropland (Percentage of total cropland)}
\begin{frame}{Top 10 Regions in Ghana with the highest acidic cropland (Percentage of total cropland)}
    \begin{figure}
    \includegraphics[width=0.95\textwidth]{D:/Dropbox/CIMMYT-WORK/Analysis/bisrat-analysis-workflows/gaia_exante_bisrat/gaia-ex-ante-wow-shared/data-output//figures_and_tables/GHA//GHA_plt_acidic_cropland_perc.png}
    \end{figure}
\end{frame}

\section*{Total crop area distribution (ha)}
\subsection*{Total crop area in the top-10 most acidic Regions, not segregated by acidic cropland}
\begin{frame}{Total crop area in the top-10 most acidic Regions, not segregated by acidic cropland}
    \begin{center}
        \includegraphics[width=0.95\textwidth]{D:/Dropbox/CIMMYT-WORK/Analysis/bisrat-analysis-workflows/gaia_exante_bisrat/gaia-ex-ante-wow-shared/data-output//figures_and_tables/GHA//GHA_plt_crop_area_acidic.png}
    \end{center}
\end{frame}

\begin{frame}{Total crop area in the top-10 most acidic Regions, not segregated by acidic cropland (table)}
    \scriptsize
    \input{D:/Dropbox/CIMMYT-WORK/Analysis/bisrat-analysis-workflows/gaia_exante_bisrat/gaia-ex-ante-wow-shared/data-output//figures_and_tables/GHA//GHA_table_crop_area_acidic.tex}
\end{frame}

\section{Additional Production}
\subsection*{Potential for additional production (tone) in top 10 acidic Regions in Ghana}
\begin{frame}{Potential for additional production (tone) in top 10 acidic Regions in Ghana}
    \begin{figure}
    \includegraphics[width=0.98\textwidth, height=0.85\textheight]{D:/Dropbox/CIMMYT-WORK/Analysis/bisrat-analysis-workflows/gaia_exante_bisrat/gaia-ex-ante-wow-shared/data-output//figures_and_tables/GHA//GHA_plt_additional_production.png}
    \end{figure}
\end{frame}

\subsection*{Potential for additional production (tone) in top 10 acidic Regions in Ghana}
\begin{frame}{Potential for additional production (1000 tone) in top 10 acidic Regions in Ghana(Table) }
    \scriptsize
    \input{D:/Dropbox/CIMMYT-WORK/Analysis/bisrat-analysis-workflows/gaia_exante_bisrat/gaia-ex-ante-wow-shared/data-output//figures_and_tables/GHA//GHA_table_additional_production.tex}
\end{frame}

\subsection*{Potential for additional production value (US\$) in top 10 acidic Regions in Ghana}
\begin{frame}{Potential for additional production value (US\$) in top 10 acidic Regions in Ghana}
    \begin{figure}
    \includegraphics[width=0.98\textwidth, height=0.85\textheight]{D:/Dropbox/CIMMYT-WORK/Analysis/bisrat-analysis-workflows/gaia_exante_bisrat/gaia-ex-ante-wow-shared/data-output//figures_and_tables/GHA//GHA_plt_additional_production_value.png}
    \end{figure}
\end{frame}

\begin{frame}{Potential for additional production value (1000 US\$) in top 10 acidic Regions in Ghana(Table)}
    \scriptsize
    \input{D:/Dropbox/CIMMYT-WORK/Analysis/bisrat-analysis-workflows/gaia_exante_bisrat/gaia-ex-ante-wow-shared/data-output//figures_and_tables/GHA//GHA_table_additional_production_value.tex}
\end{frame}

\section{Yield Loss}
\subsection*{Yield loss for Maize and Cacao(\%) in Ghana}
\begin{frame}{Yield loss due to soil acidity for Maize and Cacao in Ghana}
    \begin{figure}
    \includegraphics[width=0.95\textwidth]{D:/Dropbox/CIMMYT-WORK/Analysis/bisrat-analysis-workflows/gaia_exante_bisrat/gaia-ex-ante-wow-shared/data-output//figures_and_tables/GHA//GHA_yield_loss.png}
    \end{figure}
\end{frame}

\section{Yield Response}
\subsection*{Yield response to lime application in Ghana}
\begin{frame} {Yield response to lime application in Ghana}
    \begin{figure}
        \centering
        \includegraphics[width=0.98\textwidth]{D:/Dropbox/CIMMYT-WORK/Analysis/bisrat-analysis-workflows/gaia_exante_bisrat/gaia-ex-ante-wow-shared/data-output//figures_and_tables/GHA//GHA_yield_response.png}
    \end{figure}
\end{frame}

\section{Lime Requirement}
\subsection*{Lime requirements (t) for acid soil remediation}
\begin{frame} {Lime requirements (t) for acid soil remediation in Ghana}
    \begin{figure}
        \centering
        \includegraphics[width=0.95\textwidth]{D:/Dropbox/CIMMYT-WORK/Analysis/bisrat-analysis-workflows/gaia_exante_bisrat/gaia-ex-ante-wow-shared/data-output//figures_and_tables/GHA//GHA_plt_lime_requirements.png}
    \end{figure}
\end{frame}

\subsection*{Lime requirements (t) for acid soil remediation (table)}
\begin{frame} {Lime requirements (1000 t) for acid soil remediation in Ghana (table)}
    \scriptsize
    \input{D:/Dropbox/CIMMYT-WORK/Analysis/bisrat-analysis-workflows/gaia_exante_bisrat/gaia-ex-ante-wow-shared/data-output//figures_and_tables/GHA//GHA_table_lime_requirements.tex}
\end{frame}

\subsection*{Lime rate for Maize and Cacao in Ghana}
\begin{frame} {Lime rate for Maize and Cacao in Ghana}
    \begin{figure}
        \centering
        \includegraphics[width=0.98\textwidth]{D:/Dropbox/CIMMYT-WORK/Analysis/bisrat-analysis-workflows/gaia_exante_bisrat/gaia-ex-ante-wow-shared/data-output//figures_and_tables/GHA//GHA_lime_rate.png}
    \end{figure}
\end{frame}

\begin{frame}
\frametitle{Main assumptions and parameters}
\small
\begin{itemize}

    \item Acidic cropland is defined as having acidity saturation levels above 10\% of the Effective Cation Exchange Capacity (ECEC).
    \item Crop yields are unaffected by soil acidity until a specific acidity saturation threshold Target Acidity Saturation(TAS) is reached; beyond this threshold, yields decrease linearly, eventually reaching zero at high acidity levels.
    \item The TAS that separates maximum yield from yield decline is used as a critical value for making liming recommendations.
    \item Actual yields in acid tropical soils are a fraction of the potential yields in non-acidic soils, determined by the estimated yield loss due to acidity.
    \item Economic benefits from liming are calculated by multiplying the additional yield after soil remediation by median crop prices from sub-Saharan Africa (2016-2020 FAOSTAT data). The cost of liming is calculated using estimated lime application rates and a fixed lime price of 100 US\$/ton.
    \item The benefits of liming are assessed only for the year of application due to a lack of robust, large-scale data on long-term effects in sub-Saharan Africa.
\end{itemize}
\centering
For details, visit \href{www.acidsoils.africa}{www.acidsoils.africa}.
\end{frame}

\end{document}
